\section{Complete metric transport for finite potential loops}
\label{sec:r10-complete-loop-metric-transport}

The top-degree SHT primitive contains lower-degree coefficient pairs
which need not be closed.  R9 LTF makes that failure explicit.  There
is nevertheless a finite metric-transport identity for the entire
loop.  Its proof uses the exact boundaries of those coefficient pairs,
which telescope, rather than cancelling a curvature factor or applying
the closed-pair homotopy to a nonclosed pair.

This is an identity of actual forms.  It removes the common metric
from a specified loop whenever a single actual change of potential is
available on the whole stated domain.  It does not solve the remaining
flat-metric commuting-square problem, construct unspecified overlap
witnesses, or assert that a line bundle has a global metric potential.

\subsection{The finite loop and its full coefficient pairs}

Let $W$ be an analytic open in the original setting, let $s\ge1$, and
let $R\in B^{1,1}(W)$ be closed under both full differentials.  Supply
finitely many actual scalar potentials $\Phi_i$ and put
\[
 C_i=R-D\Phi_i,\qquad
 h_k(C_i,C_j)=\sum_{a+b=k}C_i^aC_j^b,\qquad D=d'd''.
 \tag{MT1}
\]
Every $C_i$ is closed under both differentials.  Supply a finite
oriented cycle of edges $e=(i,j)$ with signed multiplicities, and for
each oriented edge an actual pluriharmonic scalar $c_e$.  Reverse
orientation changes both the multiplicity convention and the sign of
$c_e$ in the usual way.  More concretely one may simply list a closed
edge path and use its displayed edge weights.  Require
\[
 \sum_e (\Phi_j-\Phi_i)=0,\qquad
 \sum_e c_e=0,
 \tag{MT2}
\]
and require endpoint telescoping for every function of a vertex;
equivalently the signed sum of oriented endpoints is zero.  Write
$v_e=\Phi_j-\Phi_i$ and $\overline\Phi_e=(\Phi_i+\Phi_j)/2$.

For an actual scalar $u$ and an even diagonal-degree form $Q$, set
\[
 \mathcal L(u,Q)=(u\,d''Q,\ d'u\wedge Q).
 \tag{MT3}
\]
All occurrences below have scalar first argument; expressions with
an even first argument are written componentwise when needed.  The
full two-colour identity is
\[
 \partial_1\mathcal L(u,Q)=uDQ-Du\wedge Q.
 \tag{MT4}
\]
The pair and triple spaces and the operators $\partial_1,\partial_2$
are precisely R9 CPT1.  For $k\ge0$ define actual diagonal forms
\[
 \Omega_k=\sum_e v_eh_k(C_i,C_j),\qquad
 \Gamma_k=\frac12\sum_e c_e(C_i^k+C_j^k),\qquad
 V_k=\Gamma_{k-1}-\Omega_{k-1}\quad(k\ge1).
 \tag{MT5}
\]
In particular $V_1=\Gamma_0=\Omega_0=0$ by (MT2).  For $k\ge1$
define an actual pair of degree $L_1^k$ by
\[
 f_k=-\sum_e\left\{
 \mathcal L\bigl(\overline\Phi_e,
                         v_eh_{k-1}(C_i,C_j)\bigr)
 +\mathcal L\left(c_e,
             \frac{\Phi_iC_i^{k-1}+\Phi_jC_j^{k-1}}2\right)
                      \right\}.
 \tag{MT6}
\]
The entire SHT loop of degree $s$, including its wall terms, is
\[
 \begin{split}
 P_s&=\sum_{k=1}^s f_kR^{s-k}+W_s,\\
 W_s&=-\frac12\sum_e
                 \mathcal L\bigl(|c_e|,v_eh_{s-1}(C_i,C_j)\bigr).
 \end{split}
 \tag{MT7}
\]
Thus (MT7) is the literal sum of R9 SHT6 on the finite oriented loop,
not a replacement primitive convention.  The absolute-value scalars
are actual for the unit weights used in SHT; more generally their
actual availability is part of the finite data.  No ordering chamber
is selected and no wall derivative is suppressed.

\begin{lemma}[The coefficient boundary chain]
For every $k\ge1$,
\[
                 \partial_1f_k=V_{k+1}-RV_k.
 \tag{MT8}
\]
Consequently these pairs are not asserted to be separately closed.
\end{lemma}

\begin{proof}
Since $D\Phi_i=R-C_i$ and $Dc_e=0$, (MT4) gives, edge by edge,
\[
 \begin{split}
 \partial_1\{-\mathcal L(\overline\Phi_e,v_eh_{k-1})\}
    &=Rv_eh_{k-1}
             +(\Phi_iC_j-\Phi_jC_i)h_{k-1},\\
 \partial_1\left\{-\mathcal L\left(c_e,
             \frac{\Phi_iC_i^{k-1}+\Phi_jC_j^{k-1}}2\right)\right\}
    &=-\frac{Rc_e}{2}(C_i^{k-1}+C_j^{k-1})
                        +\frac{c_e}{2}(C_i^k+C_j^k).
 \end{split}
 \tag{MT9}
\]
Here $h_{k-1}=h_{k-1}(C_i,C_j)$.  The finite polynomial identity
\[
 (\Phi_iC_j-\Phi_jC_i)h_{k-1}
       +v_eh_k(C_i,C_j)=\Phi_jC_j^k-\Phi_iC_i^k
 \tag{MT10}
\]
follows by multiplying out $h_k$.  Endpoint telescoping kills the
sum of the right side.  Summing (MT9) therefore gives
$R(\Omega_{k-1}-\Gamma_{k-1})-\Omega_k+\Gamma_k$, which is (MT8).
No actual curvature factor has been cancelled.
\end{proof}

\subsection{The homotopy with its nonclosed remainder}

For any actual scalar $h$ and any actual pair $p=(p',p'')\in L_1^k(W)$,
define the same finite triple as in R9 CPT2,
\[
 H_h(p)=\bigl(d''h\wedge p',\ -d'h\wedge p'',\ h\,d'p'\bigr).
 \tag{MT11}
\]
It lies in $L_2^{k+1}(W)$, including the negative-degree conventions
at the endpoints.  Without any closedness assumption on $p$,
\[
          \partial_2H_h(p)=Dh\wedge p+
                                  \mathcal L(h,\partial_1p).
 \tag{MT12}
\]
Indeed, writing $w=\partial_1p$, the first component is
$Dh\wedge p'-hDp'$.  Since $d''w=-Dp'$, this is
$Dh\wedge p'+h\,d''w$.  The second component is
$Dh\wedge p''+d'h\wedge(d''p''+d'p')
 =Dh\wedge p''+d'h\wedge w$.  These are exactly the two
components of (MT12).  R9 CPT3 is recovered when $w=0$.
The retained remainder in (MT12) is essential below.

\subsection{The full finite transport formula}

Choose one actual scalar $h$ on the whole $W$ and make the simultaneous
change
\[
          R'=R+Dh,\qquad \Phi'_i=\Phi_i+h.
 \tag{MT13}
\]
Then every projective curvature remains the same:
$R'-D\Phi'_i=C_i$.  Keep every edge weight $c_e$ unchanged.  Denote
the resulting literal SHT loop by $P'_s$.  Define
\[
 S_t(R',R)=\sum_{\ell=0}^t(R')^\ell R^{t-\ell}\quad(t\ge0),
 \tag{MT14}
\]
and set a sum with upper index less than its lower index equal to zero.

\begin{theorem}[Finite actual metric transport of the whole loop]
The actual triple
\[
 \boxed{\quad
 Q_{s,h}=\sum_{k=1}^{s-1}S_{s-k-1}(R',R)H_h(f_k)
 -\left(0,0,
          h\sum_{k=1}^s\Gamma_{k-1}(R')^{s-k}\right)
 \quad}
 \tag{MT15}
\]
belongs to $L_2^s(W)$ and satisfies the literal full identity
\[
                         P'_s-P_s=\partial_2Q_{s,h}.
 \tag{MT16}
\]
The statement holds for every $s\ge1$ and for the entire finite
presentation, including all valuation walls and mixed metric jets.
It does not require the individual $f_k$ to be closed.
\end{theorem}

\begin{proof}
The $v_e,C_i,\Omega_k,\Gamma_k$ are unchanged under (MT13), and
therefore the entire wall term $W_s$ is unchanged.  Expansion of (MT6)
gives the exact pair identity
\[
 f'_k-f_k=\mathcal L(h,V_k)
                         -\partial_2(0,0,h\Gamma_{k-1}).
 \tag{MT17}
\]
For example, the raw change is
$-\mathcal L(h,\Omega_{k-1})
 -(\Gamma_{k-1}d''h,h\,d'\Gamma_{k-1})$.
The identity
\[
 (\Gamma d''h,h\,d'\Gamma)+\mathcal L(h,\Gamma)
                      =\partial_2(0,0,h\Gamma)
 \tag{MT18}
\]
for any even diagonal form $\Gamma$ proves (MT17), with all signs
and derivatives retained.

Because $R,R'$ are closed, the scalar curvature telescope is
\[
             (R')^t-R^t=Dh\wedge S_{t-1}(R',R)\quad(t\ge1).
 \tag{MT19}
\]
Substitute (MT17)--(MT19) in (MT7), and apply (MT12) to each $f_k$.
The terms which are immediately $\partial_2$ boundaries give
exactly (MT15).  The remaining pair is
\[
 \begin{split}
 &\sum_{k=1}^s(R')^{s-k}\mathcal L(h,V_k)\\
 &\hspace{7mm}-\sum_{k=1}^{s-1}S_{s-k-1}(R',R)
                                      \mathcal L(h,V_{k+1}-RV_k).
 \end{split}
 \tag{MT20}
\]
Since $R$ is closed, multiplication by $R$ commutes with
$\mathcal L(h,\cdot)$.  The coefficient of
$\mathcal L(h,V_k)$ for $2\le k\le s-1$ is
\[
              (R')^{s-k}+R S_{s-k-1}-S_{s-k}=0.
 \tag{MT21}
\]
The same cancellation at $k=s$ reads $1-S_0=0$.
The possible $k=1$ endpoint vanishes because $V_1=0$.
Thus (MT20) is zero.  This proves (MT16).  At $s=1$ the first
sum in (MT15) is empty and $\Gamma_0=0$, so the complete loop
is literally invariant under the common change; this also follows
directly from SHT6.
\end{proof}

One can alternatively set $R_t=R+tDh$, $\Phi_i(t)=\Phi_i+th$.
The resulting $P_s(t)$ is a finite polynomial in the numerical
parameter $t$ with actual coefficients.  Its derivative is
\[
 \frac{dP_s(t)}{dt}=\partial_2\left\{
 \sum_{k=1}^{s-1}(s-k)R_t^{s-k-1}H_h(f_k(t))
 -\left(0,0,h\sum_{k=1}^s\Gamma_{k-1}R_t^{s-k}\right)
                                  \right\}.
 \tag{MT22}
\]
Integration in this formula means only integrating the finitely many
real polynomial coefficients.  It introduces no actual-form
completion, analytic pushforward, or parameterwise choice of a
source layer.  The finite formula (MT15) already suffices.

\subsection{Exact scope and the flat-metric reduction}

For the two commuting distinct-entry changes in R9 PPD, take the
oriented loop $0\to A\to T\to B\to0$ with weights
$a,b,-a,-b$.  Then $P_s$ is exactly the PPD8 full chosen-primitive
defect.  The theorem gives its complete common-metric transport,
including all four triangle, absolute-wall, and metric-transport
contributions of PPD9.  It does not replace that residual by its
projective-fibre period.

If on one whole neighborhood $N$ of the cancelled set the line bundle
has a specified analytic frame and its chosen metric has an actual
potential $h$, let $P_s^{\rm flat}$ be the same full loop computed
with the norm-one frame metric.  On $W=N\setminus|f=0|$ the theorem
constructs $Q_{s,h}$ with
\[
                  P_s^{\rm metric}-P_s^{\rm flat}
                                            =\partial_2Q_{s,h}.
 \tag{MT23}
\]
Thus any actual relative witness
$P_s^{\rm flat}=\partial_2q+e|_W$, with $e$ defined on all of $N$,
transports explicitly to the original metric by replacing $q$ with
$q+Q_{s,h}$.  Conversely subtracting $Q_{s,h}$ transports a supplied
witness in the other direction.  This is a reduction of an overlap
problem; it is not a proof that the flat-metric witness exists.

When the expected extending pair is
$e_{a,b}R^{s-1}$, with
$e_{a,b}=(a\,d''b,b\,d'a)$, its change under a supplied common
metric potential is also an explicit $\partial_2$ boundary by
R9 CPT, since $e_{a,b}$ is independently closed.  Hence the full-loop
metric transport is consistent with the ASP period normalization
and does not attempt to remove that genuine extending-pair term.

A cover by frame domains is not one global $h$.  Applying (MT23)
on each frame domain does not by itself glue the resulting triples;
the separate RAG whole-overlap hypotheses remain necessary.  Nor
does the theorem supply a finite elementary factorization of an
arbitrary analytic matrix.  In particular the general distinct-entry
flat-metric square, arbitrary frame-path coherence, full lower-middle
OPEN-KERNEL, general strong G, arbitrary higher actual descent and
complete Beilinson--Bloch height theory remain open unless settled
by separate arguments.

Every expression above uses finite actual addition, multiplication,
full differentiation and one fixed finite collection of scalar
potentials.  It commutes with actual restriction and refinement.
The pair and triple spaces carry their own inherited coefficient-final
topologies.  On one fixed whole coefficient presentation all formulas
retain the complete mixed jets, and the finite Leibniz estimates give
finite-stage continuity into one common finite refinement.  No global
wedge-continuity assertion, Hausdorff quotient, strict image theorem,
or comparison with a probe topology is inferred.

\paragraph{Status and provenance.}
The finite identity (MT16), its nonclosed remainder calculation and
the conditional reduction (MT23) are proved in R10.  The actual SHT
primitive convention, CPT pair complex, and LTF failure are frozen
R9 inputs at their indexed hashes.  The root and the diagonal
constructor independently found the same telescope and are both
construction contributors; neither is an independent reviewer of
this module.  No literature-novelty claim is made.
